
\chapter{Introduction}

\section{Motivation}
Automated program repair (APR) is the task of locating and fixing the defective behaviour in code.
Success is typically defined as whether the change proposed by an APR system resolves this faulty behaviour. 
Formally, given a task description $x$ and a buggy program $\tilde{y}$, APR aims to produce a fixed program $\hat{y}$
such that $\hat{y}$ passes all test cases $t \in T$ in an executor $\mathcal{E}$, i.e., $\mathcal{E}(\hat{y}, T) = 1$ while $\mathcal{E}(\tilde{y}, T) = 0$.

Early APR approaches mostly relied on symbolic analysis, such as abstract syntax trees (ASTs), predefined fix templates, 
search-based methods that mutated code until tests passed, or constraint solvers that synthesized repairs based on formal specifications.
Modern systems have shifted to probabilistic analysis with flexible reasoning based on non-deterministic models that are trained on large code and text corpora.
Recent surveys show that agent-based repair becomes a dominant paradigm~\citep{zhang2024slrapr, liu2024agentsse}.
Here, agents are understood as iterative processes that are built around an underlying LLM model and repeatedly invoke it while keeping a context.
Agents often interact with the environment through requesting files content, search for symbols, and run tests.
This loop iterates until a desired goal is achieved.

Most of current agent systems share a common assumption that static source code is enough for repair,
although many systems provide a high-level access to the code.
For example, \citet{zhang2024autocoderover} proposed AutoCodeRover, which replaces raw file reads with AST-level retrieval and
spectrum-based fault localization. \citet{xia2024agentless} designed
Agentless, a fixed pipeline of hierarchical localization, patch sampling, and test-based candidate selection.
On SWE-bench Lite \citep{jimenez2023swebench}, these three systems reached 12.5\%, 19\%, and 32\% pass@1, respectively,
where pass@1 measures the proportion of test samples resolved correctly by the model on the first try.

However, a fundamental limitation of this approach is that this is not how humans actually debug programs.
When investigating a failed test, a developer can break the program in a debugger, read the call stack, inspect how variables
change over time, and check the program state at the crash point.
Systems that perform APR based on just static source code do not give the LLM model any of the runtime information that human
relies on when fixing a bug.
Instead, the only available information is the source code and an exception message,
and the model must reconstruct what the program did from static context alone.
Even worse, sometimes the only available information is incorrect program output, which makes bug localization even more challenging.
To demonstrate the importance of accessing runtime program state, consider the following classical example of off-by-one error:

\begin{algorithm}[ht]
\caption{Binary search with an off-by-one bug: \texttt{hi} is initialized to \texttt{len(nums)} instead of \texttt{len(nums) - 1}.}
\label{alg:binary-search-bug}
\begin{algorithmic}[1]
\Require sorted array \texttt{nums}, integer \texttt{target}
\Ensure index of \texttt{target} in \texttt{nums}, or $-1$ if absent
\State $lo \gets 0$
\State $hi \gets \text{len}(nums)$
\While{$lo < hi$}
    \State $mid \gets \lfloor (lo + hi) / 2 \rfloor$
    \If{$nums[mid] = target$}
        \Return $mid$
    \ElsIf{$nums[mid] < target$}
        \State $lo \gets mid + 1$
    \Else
        \State $hi \gets mid$
    \EndIf
\EndWhile
\State \Return $-1$
\end{algorithmic}
\end{algorithm}

While it is possible to find this subtle bug by carefully verifying the algorithm, the reader can imagine how complicated
real-world implementations can get, where such bugs can become even more elusive.

This thesis is not the first one to point out this limitation. \citet{chen2023teaching} point to this as the core constraint of
execution-free repair, and \citet{zhong2024debug} show that explicitly providing models
with control flow and variable state information before patch generation boosts performance.

Our goal is to investigate whether using runtime information, such as
the call stack, value of local variables, and execution trace, for program repair
improves patch correctness. We evaluate this approach on a baseline with access to only static source code with
bugs in real-world Python projects.
We build on prior evidence that such execution-derived context helps model reasoning on coding tasks
\citep{ni2024nextteachinglargelanguage,han2025rethinkingcapabilityfinetunedlanguage}.

\section{Our contributions}
\begin{enumerate}
  \item We develop a custom debugger hooks for trace collection which 
      allow collecting runtime frames in the most common Python test frameworks: \texttt{pytest}, \texttt{unittest}, \texttt{django} (modification of \texttt{unittest}).
      The collected traces are serialized in a structured JSON format and contain entry/exit frames
        and line-by-line frames with the values of local variables.
  \item We perform comparative and ablative experiments. On SWE-bench Lite and BugsInPy, we
    measure patch correctness in an interactive environment, where model can perform source code exploration and, optionally,
    has access to collected runtime information.
    We evaluate under 2 settings: baseline where model has access to no runtime information, and where runtime information is available.
    Ablations over trace subsets quantify how much each category of runtime information contributes to repair success.
  \item We fine-tune a small model on the collected runtime data to explore how runtime information changes model performance.
      Specifically, we run experiments on Qwen2.5-Coder-7B-Instruct on approximately 300 training samples
      collected using our custom debugger hooks. We train with QLoRA on a single A100 GPU across five seeds.
      We treat fine-tuning as an open empirical question and focus primarily on inference-time use of runtime information.
\end{enumerate}

\section{Thesis structure}
This thesis attempts to address the following research questions:

\begin{tcolorbox}[colback=blue!6!white,colframe=blue!60!black,boxrule=0.6pt,arc=2pt,left=6pt,right=6pt,top=4pt,bottom=4pt]
\textbf{RQ1:} Can execution traces improve LLM-based bug repairing compared to static analysis alone?
\end{tcolorbox}

\begin{tcolorbox}[colback=green!6!white,colframe=green!50!black,boxrule=0.6pt,arc=2pt,left=6pt,right=6pt,top=4pt,bottom=4pt]
\textbf{RQ2:} What is the minimal trace information needed for effective debugging assistance?
\end{tcolorbox}

\begin{tcolorbox}[colback=orange!8!white,colframe=orange!70!black,boxrule=0.6pt,arc=2pt,left=6pt,right=6pt,top=4pt,bottom=4pt]
\textbf{RQ3:} Does fine-tuning on debugging-specific data outperform prompting general-purpose models? 
\end{tcolorbox}

\Cref{ch:related-work} reviews the related work on modern APR, LLM-based repair, and recent approaches that involve runtime frames in resolving incorrect program behaviour.
\Cref{ch:system-methodology} presents our approach, including runtime information collection, filtering and serialization of runtime frames, and the LLM tool-based interface.
The evaluation chapter describes the experimental setup, datasets, metrics, and results, including ablations over runtime context.
\Cref{ch:discussion} discusses limitations and threats to validity.
\Cref{ch:conclusion} summarizes our results, and outlines future work and possible improvements.
\change{TODO: validate this paragraph when we finish later chapters}
